\appendix
\chapter{Annexes}

\section{Structure du dépôt de code}

\begin{lstlisting}[language={}, caption={Arborescence simplifiée du projet}]
alert_receiver.py          Client Wazuh Indexer (recherche/lecture d'alertes)
poller_service.py          Boucle de collecte periodique, watermark
investigation_policy.py    Moteur de politiques (quelles alertes ouvrir)
investigation_manager.py   Creation des investigations, liaison des actifs
database.py                Schema MySQL (DDL + migrations) + acces donnees

ai_investigator/
  investigator.py          Boucle agentique (agent d'investigation)
  chatbot.py                Assistant conversationnel (memes outils)
  llm/                      Interface LLMProvider (API / local)
  tools/                    Outils Wazuh + base de donnees + renseignement

api/
  main.py                   Application FastAPI, cycle de vie, routes
  dependencies.py            Cablage des services (DB, poller, IA, LLM)
  auth.py                    Verification des identifiants + JWT
  routers/                   auth, investigations, policies, assets,
                              users, system (poller), chat

frontend/
  src/views/                Tableau de bord, Investigations, Actifs,
                              Politiques, Parametres, Chat, Connexion
  src/stores/                Etat Pinia (investigations, chat, auth)
  src/api/                   Client HTTP + modules par ressource
\end{lstlisting}

\section{Extrait du schéma de la table \texttt{investigations}}

\begin{lstlisting}[language=SQL, caption={Extrait de définition de la table
\texttt{investigations}}]
CREATE TABLE IF NOT EXISTS investigations (
    id          CHAR(36)     PRIMARY KEY,
    alert_id    VARCHAR(255) NOT NULL,
    asset_id    CHAR(36),
    policy_id   CHAR(36),
    title       VARCHAR(255),
    description TEXT,
    status      VARCHAR(20)  NOT NULL DEFAULT 'QUEUED',
    severity    VARCHAR(20),
    closed_at   DATETIME,
    created_at  DATETIME     NOT NULL DEFAULT CURRENT_TIMESTAMP,
    updated_at  DATETIME     NOT NULL DEFAULT CURRENT_TIMESTAMP
                ON UPDATE CURRENT_TIMESTAMP,
    FOREIGN KEY (asset_id)  REFERENCES assets(id),
    FOREIGN KEY (policy_id) REFERENCES investigation_policies(id)
);
\end{lstlisting}

\section{Extrait de l'invite système de l'agent d'investigation}

\begin{lstlisting}[language={}, caption={Extrait du \textit{system prompt}
de l'AI Investigator}]
You are an expert AI security analyst working on a Security
Investigation Platform.

Task: investigate a security alert and determine TRUE POSITIVE,
FALSE POSITIVE, or NEEDS REVIEW.

## Process
1. Read the triggering alert AND the baseline context below it
   (asset record + recent alert history for the same agent).
2. Only call tools for what's still missing.
3. Correlate the evidence, then give your verdict.

## Be efficient - budget matters as much as accuracy
- Reach a verdict in 2-4 tool calls for most alerts.
- Duplicate calls (same tool + same arguments) are blocked.
- ...
\end{lstlisting}

\section{Variables de configuration principales (\texttt{.env})}

\begin{table}[H]
\centering
\caption{Principales variables d'environnement}
\begin{tabular}{@{}p{5.5cm}p{8cm}@{}}
\toprule
\textbf{Variable} & \textbf{Rôle} \\
\midrule
\texttt{WAZUH\_INDEXER\_HOST/PORT/\allowbreak USERNAME/PASSWORD} & Connexion
au Wazuh Indexer (compte de service applicatif) \\
\texttt{MYSQL\_HOST/PORT/USER/\allowbreak PASSWORD/DATABASE} & Connexion à
la base applicative \\
\texttt{LLM\_PROVIDER} & \texttt{local} (Ollama) ou \texttt{api} (tout
point de terminaison compatible OpenAI) \\
\texttt{LLM\_MODEL}, \texttt{LLM\_API\_KEY} & Modèle utilisé et clé
d'accès le cas échéant \\
\texttt{LLM\_MAX\_TOKENS}, \texttt{LLM\_REASONING\_EFFORT} & Bornes de
génération et effort de raisonnement (selon le modèle) \\
\texttt{ABUSEIPDB\_API\_KEY} & Active l'outil \texttt{check\_ip\_reputation}
si renseignée \\
\texttt{AUTH\_JWT\_SECRET}, \texttt{AUTH\_TOKEN\_TTL\_MINUTES} & Signature
et durée de vie des jetons de session \\
\bottomrule
\end{tabular}
\end{table}

\section{Glossaire complémentaire}

\begin{description}[leftmargin=3.4cm]
  \item[Agent (Wazuh)] Logiciel installé sur un hôte surveillé, chargé de
    collecter les événements et de les transmettre au Wazuh Manager.
  \item[Alerte] Événement de sécurité généré lorsqu'une règle de détection
    est déclenchée par un événement collecté.
  \item[Boucle agentique] Cycle itératif dans lequel un modèle de langage
    alterne raisonnement et appels d'outils jusqu'à produire une conclusion.
  \item[Faux positif] Alerte ne correspondant à aucune menace réelle.
  \item[Verdict] Conclusion produite par l'agent d'investigation
    (\texttt{TRUE\_POSITIVE}, \texttt{FALSE\_POSITIVE},
    \texttt{NEEDS\_REVIEW}), avec un niveau de confiance et une
    justification.
\end{description}
